% Seminar 3 - Pagina 10. Fragment pentru
% % Seminar 3 - Pagina 10. Fragment pentru
% % Seminar 3 - Pagina 10. Fragment pentru
% % Seminar 3 - Pagina 10. Fragment pentru
% \input{seminars/pages/seminar3-pagina10-ro}; fără preambul sau pagină nouă.
% NOTE EDITORIALE:
% - Această ultimă pagină continuă discuția despre cardinalitate de pe
%   pagina 9 a Seminarului 3.
% - Pentru cardinalul 1, mulțimea manuscrisului pare a fi {{vid}},
%   nu {vid}. Ambele au un singur element, deci au cardinalul 1.
% - La suma cardinalelor, manuscrisul pare să scrie A_i REUNIUNE {i},
%   nu A_i PRODUS CARTEZIAN {i}. Prima formulă nu dă în general o
%   reuniune disjunctă și nici suma cardinalelor. Formula standard este
%     | \bigcup_{i\in I} (A_i \times \{i\}) |.
%   Am păstrat citirea aparentă a manuscrisului în formula afișată.
% - Manuscrisul afirmă că definițiile sunt corecte fără a demonstra că
%   ele nu depind de alegerea mulțimilor reprezentante A_i, A și B.
% - "Paradoxul lui Cantor" este redat ca în notițe. Distincția modernă
%   dintre mulțime și clasă proprie nu este adăugată în text.
% - Ultimele comparații de cardinali nu dezvoltă premisele axiomatice;
%   ele sunt păstrate în forma din manuscris.

\subsection*{Seminar 3 - Pagina 10}
\begingroup
\setlength{\abovedisplayskip}{5pt}
\setlength{\belowdisplayskip}{5pt}
\setlength{\abovedisplayshortskip}{3pt}
\setlength{\belowdisplayshortskip}{3pt}

\noindent Primele numere cardinale sunt notate astfel:
\[
\begin{aligned}
0 &:= \operatorname{card}(\varnothing),\\
1 &:= \operatorname{card}\bigl(\{\{\varnothing\}\}\bigr),\\
2 &:= \operatorname{card}\bigl(\{\{\varnothing\},\{\{\varnothing\}\}\}\bigr).
\end{aligned}
\]

\medskip
\noindent\underline{Operații cu numere cardinale.}

Fie $\alpha=|A|$. Dacă $(\alpha_i)_{i\in I}$ este o familie de numere
cardinale, cu $\alpha_i=|A_i|$, definim
\[
\sum_{i\in I}\alpha_i\ \overset{\mathrm{def}}{=}
\left|\bigcup_{i\in I}\bigl(A_i\cup\{i\}\bigr)\right|.
\]

\noindent Pentru $\alpha=|A|$ și $\beta=|B|$, definim
\[
\alpha\cdot\beta\ \overset{\mathrm{def}}{=}\ |A\times B|,
\qquad
\alpha^{\beta}\ \overset{\mathrm{def}}{=}
\bigl|\operatorname{Hom}(B,A)\bigr|.
\]

\noindent Definițiile sunt corecte.

\medskip
\noindent\underline{Paradoxul lui Cantor.}
Mulțimea tuturor numerelor cardinale este contradictorie
(se obține $\alpha<\alpha$).

\medskip
\noindent Pentru $\alpha=|A|$ și $\beta=|B|$,
\[
\alpha\leq\beta\ \iff
\exists\, f:A\longrightarrow B\text{ injectivă}.
\]
Inegalitatea strictă este dată prin
\[
\alpha<\beta\quad\iff\quad
\begin{cases}
\text{există o aplicație injectivă }f:A\longrightarrow B,\\
\text{nu există o aplicație bijectivă }A\longrightarrow B.
\end{cases}
\]

\noindent Definiția este corectă.
\endgroup
; fără preambul sau pagină nouă.
% NOTE EDITORIALE:
% - Această ultimă pagină continuă discuția despre cardinalitate de pe
%   pagina 9 a Seminarului 3.
% - Pentru cardinalul 1, mulțimea manuscrisului pare a fi {{vid}},
%   nu {vid}. Ambele au un singur element, deci au cardinalul 1.
% - La suma cardinalelor, manuscrisul pare să scrie A_i REUNIUNE {i},
%   nu A_i PRODUS CARTEZIAN {i}. Prima formulă nu dă în general o
%   reuniune disjunctă și nici suma cardinalelor. Formula standard este
%     | \bigcup_{i\in I} (A_i \times \{i\}) |.
%   Am păstrat citirea aparentă a manuscrisului în formula afișată.
% - Manuscrisul afirmă că definițiile sunt corecte fără a demonstra că
%   ele nu depind de alegerea mulțimilor reprezentante A_i, A și B.
% - "Paradoxul lui Cantor" este redat ca în notițe. Distincția modernă
%   dintre mulțime și clasă proprie nu este adăugată în text.
% - Ultimele comparații de cardinali nu dezvoltă premisele axiomatice;
%   ele sunt păstrate în forma din manuscris.

\subsection*{Seminar 3 - Pagina 10}
\begingroup
\setlength{\abovedisplayskip}{5pt}
\setlength{\belowdisplayskip}{5pt}
\setlength{\abovedisplayshortskip}{3pt}
\setlength{\belowdisplayshortskip}{3pt}

\noindent Primele numere cardinale sunt notate astfel:
\[
\begin{aligned}
0 &:= \operatorname{card}(\varnothing),\\
1 &:= \operatorname{card}\bigl(\{\{\varnothing\}\}\bigr),\\
2 &:= \operatorname{card}\bigl(\{\{\varnothing\},\{\{\varnothing\}\}\}\bigr).
\end{aligned}
\]

\medskip
\noindent\underline{Operații cu numere cardinale.}

Fie $\alpha=|A|$. Dacă $(\alpha_i)_{i\in I}$ este o familie de numere
cardinale, cu $\alpha_i=|A_i|$, definim
\[
\sum_{i\in I}\alpha_i\ \overset{\mathrm{def}}{=}
\left|\bigcup_{i\in I}\bigl(A_i\cup\{i\}\bigr)\right|.
\]

\noindent Pentru $\alpha=|A|$ și $\beta=|B|$, definim
\[
\alpha\cdot\beta\ \overset{\mathrm{def}}{=}\ |A\times B|,
\qquad
\alpha^{\beta}\ \overset{\mathrm{def}}{=}
\bigl|\operatorname{Hom}(B,A)\bigr|.
\]

\noindent Definițiile sunt corecte.

\medskip
\noindent\underline{Paradoxul lui Cantor.}
Mulțimea tuturor numerelor cardinale este contradictorie
(se obține $\alpha<\alpha$).

\medskip
\noindent Pentru $\alpha=|A|$ și $\beta=|B|$,
\[
\alpha\leq\beta\ \iff
\exists\, f:A\longrightarrow B\text{ injectivă}.
\]
Inegalitatea strictă este dată prin
\[
\alpha<\beta\quad\iff\quad
\begin{cases}
\text{există o aplicație injectivă }f:A\longrightarrow B,\\
\text{nu există o aplicație bijectivă }A\longrightarrow B.
\end{cases}
\]

\noindent Definiția este corectă.
\endgroup
; fără preambul sau pagină nouă.
% NOTE EDITORIALE:
% - Această ultimă pagină continuă discuția despre cardinalitate de pe
%   pagina 9 a Seminarului 3.
% - Pentru cardinalul 1, mulțimea manuscrisului pare a fi {{vid}},
%   nu {vid}. Ambele au un singur element, deci au cardinalul 1.
% - La suma cardinalelor, manuscrisul pare să scrie A_i REUNIUNE {i},
%   nu A_i PRODUS CARTEZIAN {i}. Prima formulă nu dă în general o
%   reuniune disjunctă și nici suma cardinalelor. Formula standard este
%     | \bigcup_{i\in I} (A_i \times \{i\}) |.
%   Am păstrat citirea aparentă a manuscrisului în formula afișată.
% - Manuscrisul afirmă că definițiile sunt corecte fără a demonstra că
%   ele nu depind de alegerea mulțimilor reprezentante A_i, A și B.
% - "Paradoxul lui Cantor" este redat ca în notițe. Distincția modernă
%   dintre mulțime și clasă proprie nu este adăugată în text.
% - Ultimele comparații de cardinali nu dezvoltă premisele axiomatice;
%   ele sunt păstrate în forma din manuscris.

\subsection*{Seminar 3 - Pagina 10}
\begingroup
\setlength{\abovedisplayskip}{5pt}
\setlength{\belowdisplayskip}{5pt}
\setlength{\abovedisplayshortskip}{3pt}
\setlength{\belowdisplayshortskip}{3pt}

\noindent Primele numere cardinale sunt notate astfel:
\[
\begin{aligned}
0 &:= \operatorname{card}(\varnothing),\\
1 &:= \operatorname{card}\bigl(\{\{\varnothing\}\}\bigr),\\
2 &:= \operatorname{card}\bigl(\{\{\varnothing\},\{\{\varnothing\}\}\}\bigr).
\end{aligned}
\]

\medskip
\noindent\underline{Operații cu numere cardinale.}

Fie $\alpha=|A|$. Dacă $(\alpha_i)_{i\in I}$ este o familie de numere
cardinale, cu $\alpha_i=|A_i|$, definim
\[
\sum_{i\in I}\alpha_i\ \overset{\mathrm{def}}{=}
\left|\bigcup_{i\in I}\bigl(A_i\cup\{i\}\bigr)\right|.
\]

\noindent Pentru $\alpha=|A|$ și $\beta=|B|$, definim
\[
\alpha\cdot\beta\ \overset{\mathrm{def}}{=}\ |A\times B|,
\qquad
\alpha^{\beta}\ \overset{\mathrm{def}}{=}
\bigl|\operatorname{Hom}(B,A)\bigr|.
\]

\noindent Definițiile sunt corecte.

\medskip
\noindent\underline{Paradoxul lui Cantor.}
Mulțimea tuturor numerelor cardinale este contradictorie
(se obține $\alpha<\alpha$).

\medskip
\noindent Pentru $\alpha=|A|$ și $\beta=|B|$,
\[
\alpha\leq\beta\ \iff
\exists\, f:A\longrightarrow B\text{ injectivă}.
\]
Inegalitatea strictă este dată prin
\[
\alpha<\beta\quad\iff\quad
\begin{cases}
\text{există o aplicație injectivă }f:A\longrightarrow B,\\
\text{nu există o aplicație bijectivă }A\longrightarrow B.
\end{cases}
\]

\noindent Definiția este corectă.
\endgroup
; fără preambul sau pagină nouă.
% NOTE EDITORIALE:
% - Această ultimă pagină continuă discuția despre cardinalitate de pe
%   pagina 9 a Seminarului 3.
% - Pentru cardinalul 1, mulțimea manuscrisului pare a fi {{vid}},
%   nu {vid}. Ambele au un singur element, deci au cardinalul 1.
% - La suma cardinalelor, manuscrisul pare să scrie A_i REUNIUNE {i},
%   nu A_i PRODUS CARTEZIAN {i}. Prima formulă nu dă în general o
%   reuniune disjunctă și nici suma cardinalelor. Formula standard este
%     | \bigcup_{i\in I} (A_i \times \{i\}) |.
%   Am păstrat citirea aparentă a manuscrisului în formula afișată.
% - Manuscrisul afirmă că definițiile sunt corecte fără a demonstra că
%   ele nu depind de alegerea mulțimilor reprezentante A_i, A și B.
% - "Paradoxul lui Cantor" este redat ca în notițe. Distincția modernă
%   dintre mulțime și clasă proprie nu este adăugată în text.
% - Ultimele comparații de cardinali nu dezvoltă premisele axiomatice;
%   ele sunt păstrate în forma din manuscris.

\subsection*{Seminar 3 - Pagina 10}
\begingroup
\setlength{\abovedisplayskip}{5pt}
\setlength{\belowdisplayskip}{5pt}
\setlength{\abovedisplayshortskip}{3pt}
\setlength{\belowdisplayshortskip}{3pt}

\noindent Primele numere cardinale sunt notate astfel:
\[
\begin{aligned}
0 &:= \operatorname{card}(\varnothing),\\
1 &:= \operatorname{card}\bigl(\{\{\varnothing\}\}\bigr),\\
2 &:= \operatorname{card}\bigl(\{\{\varnothing\},\{\{\varnothing\}\}\}\bigr).
\end{aligned}
\]

\medskip
\noindent\underline{Operații cu numere cardinale.}

Fie $\alpha=|A|$. Dacă $(\alpha_i)_{i\in I}$ este o familie de numere
cardinale, cu $\alpha_i=|A_i|$, definim
\[
\sum_{i\in I}\alpha_i\ \overset{\mathrm{def}}{=}
\left|\bigcup_{i\in I}\bigl(A_i\cup\{i\}\bigr)\right|.
\]

\noindent Pentru $\alpha=|A|$ și $\beta=|B|$, definim
\[
\alpha\cdot\beta\ \overset{\mathrm{def}}{=}\ |A\times B|,
\qquad
\alpha^{\beta}\ \overset{\mathrm{def}}{=}
\bigl|\operatorname{Hom}(B,A)\bigr|.
\]

\noindent Definițiile sunt corecte.

\medskip
\noindent\underline{Paradoxul lui Cantor.}
Mulțimea tuturor numerelor cardinale este contradictorie
(se obține $\alpha<\alpha$).

\medskip
\noindent Pentru $\alpha=|A|$ și $\beta=|B|$,
\[
\alpha\leq\beta\ \iff
\exists\, f:A\longrightarrow B\text{ injectivă}.
\]
Inegalitatea strictă este dată prin
\[
\alpha<\beta\quad\iff\quad
\begin{cases}
\text{există o aplicație injectivă }f:A\longrightarrow B,\\
\text{nu există o aplicație bijectivă }A\longrightarrow B.
\end{cases}
\]

\noindent Definiția este corectă.
\endgroup
